% SPDX-License-Identifier: Apache-2.0
\section{A Unit Between Two Channels, Lowered}
\label{sec:x-stage}

The lowering meets the handshake. A stage waits, with \code{until},
until a word is offered on its input and its consumer can take one,
then passes the word on incremented and counts it. In the lowering a
channel port is a valid/ready handshake around its data: an
\code{Rx<T>} is \code{data} and \code{valid} in and \code{ready} out,
a \code{Tx<T>} the reverse. The wait becomes the guard on every
register drive after it; a receive asserts \code{ready} under the
guard; a send drives \code{data} and asserts \code{valid} under it.
What comes out, Listing~\ref{lst:x-stage-out}, is a combinational
stage with a handshake and one register, which is what the source
says, and Figure~\ref{fig:x-stage} is the same stage simulated, with
the source offering every other cycle.

\begin{figure}[htbp]
\centering
\input{stage_timing.tex}
\caption{The stage: a word in every other cycle, the same word plus one out in the same cycle, the count of words passed.}
\label{fig:x-stage}
\end{figure}

\lstinputlisting[caption={\code{ex\_stage.rs}. Run with \code{bazel run //lib/examples:ex\_stage}.},label={lst:x-stage}]{lib/examples/ex_stage.rs}

\lstinputlisting[style=ascii,caption={What \code{ex\_stage} printed when the build ran it: the simulation, then the lowering.},label={lst:x-stage-out}]{out_stage.txt}
